% ============================================================
% The Alignment Limit for All Primes
% ============================================================

\documentclass[12pt]{amsart}

\usepackage{amsmath,amssymb,amsthm}
\usepackage{booktabs}
\usepackage{microtype}
\usepackage{url}

\newtheorem{theorem}{Theorem}
\newtheorem{lemma}[theorem]{Lemma}
\newtheorem{proposition}[theorem]{Proposition}
\newtheorem{corollary}[theorem]{Corollary}
\newtheorem{conjecture}[theorem]{Conjecture}
\theoremstyle{definition}
\newtheorem{definition}[theorem]{Definition}
\theoremstyle{remark}
\newtheorem{remark}[theorem]{Remark}

\title{The Alignment Limit for All Primes}

\author{Alexander S.\ Petty}

\email{alexander.petty@gmail.com}
\date{August 2021 (revised April 2026)}
\subjclass[2020]{11A63, 11B83}

\begin{document}
\begin{abstract}
The companion papers proved that the repetend alignment of
$n = pm$ (with $m$ a $b$-smooth integer and $p$ a
digit-partitioning prime) equals $(2m{-}1)/(pm{-}1)$.
This paper gives an exact formula for all primes.

For any prime $p$ not dividing $b$ and any $b$-smooth
$m \ge 1$:
\[
\alpha(pm) = \frac{m - 1 + m \cdot T(C_m)/L}{pm - 1},
\]
where $L = \operatorname{ord}_p(b)$, the coset
$C_m = u\langle b \rangle$ is determined by $u = b^t/m$
with $m \mid b^t$, and $T(C_m) = \sum_{r \in C_m}
n_{\delta(r)}$ is the coset bin sum. The formula is
exact at every finite $m$.

The subsequential alignment limit as $m \to \infty$
depends on the coset $C_m$. When $b$ is a primitive
root modulo $p$, there is one coset and the limit
is $(p{-}1+S)/(p(p{-}1))$ where $S = \sum n_d^2$.
In general, different cosets yield different limits.

For all primes $p \ge 5$, all cosets, and all
$b$-smooth $m$: $\alpha(pm) < 1/\varphi$. Among
primes $p \ge 3$, only $p = 3$ produces alignment
above the golden threshold.
\end{abstract}
\maketitle

% ============================================================
\section{Introduction}

This paper extends the alignment formula of the preceding papers to all primes.

The first paper in this series~\cite{paper1} proved that the
repetend alignment of $n = 3m$ equals $(2m{-}1)/(3m{-}1)$ and that
the golden ratio's minimal polynomial uniquely divides the
self-referential cubic at $p = 3$. The second paper~\cite{paper2} extended the alignment
formula to all digit-partitioning primes ($p \le b{+}1$) and
characterized them by the injectivity of the digit function.

For primes $p > b{+}1$, the digit function is not
injective. Fractions from different residue classes share
digits at common positions, and the alignment formula
$(2m{-}1)/(pm{-}1)$ underestimates the true alignment.

This paper computes the correction. We show that the alignment
exact formula at every finite $m$ is determined by the
coset bin sum. The alignment never reaches $1/\varphi$
for any prime $p \ge 5$.

% ============================================================
\section{The Digit Function and Its Bins}

\begin{definition}
For a prime $p$ not dividing $b$, the \emph{digit function} is
$\delta(r) = \lfloor br/p \rfloor$ for $r \in \{0, \ldots, p{-}1\}$.
The \emph{bin} of digit $d$ is
$B_d = \{r \in \{1, \ldots, p{-}1\} : \delta(r) = d\}$, and its
size is $n_d = |B_d|$.
\end{definition}

The $p - 1$ nonzero remainders are distributed among $b$ bins. Let
$q = \lfloor (p{-}1)/b \rfloor$ and $r = (p{-}1) \bmod b$. Then $r$
bins have size $q + 1$ and $b - r$ bins have size $q$.

\begin{lemma}\label{lem:bin-sum}
The sum of squared bin sizes is
\[
S(p,b) = \sum_{d=0}^{b-1} n_d^2 = bq^2 + r(2q + 1).
\]
\end{lemma}

\begin{proof}
There are $r$ bins of size $q{+}1$ and $b{-}r$ bins of size $q$.
Thus
\[
S = r(q{+}1)^2 + (b{-}r)q^2 = rq^2 + 2rq + r + bq^2 - rq^2
  = bq^2 + r(2q + 1). \qedhere
\]
\end{proof}

\begin{corollary}\label{cor:dp-recovery}
When $p \le b + 1$, all bins have size at most $1$, so
$S(p,b) = p - 1$.
\end{corollary}

% ============================================================
\section{The Coset and Its Bin Sum}

For $n = pm$ with $m$ a $b$-smooth integer, choose $t$
with $m \mid b^t$ and set $u = b^t / m$. The coset
$C_m = u \langle b \rangle \bmod p$ is independent of the
choice of~$t$: if $m \mid b^{t'}$ as well, then
$b^{t'}/m = (b^{t}/m) \cdot b^{t'-t}$, which lies in the
same coset of $\langle b \rangle$.
After $t$ prefix digits, the repetend of $k/(pm)$ is the
repetend of $uk/p$~\cite{paper2}. The relevant orbit is
not the subgroup $\langle b \rangle$ but the coset~$C_m$.

\begin{definition}
For a coset $C$ of $\langle b \rangle$ in
$(\mathbb{Z}/p\mathbb{Z})^{*}$, the \emph{coset bin sum}
is
\[
T(C) = \sum_{r \in C} n_{\delta(r)},
\]
where $n_d = |B_d|$ is the size of the $d$-th bin.
\end{definition}

Different $b$-smooth values of $m$ can yield different
cosets and therefore different bin sums.

\begin{remark}\label{rem:equidist}
The coset $C$ has \emph{equidistributed bin sum} if
$T(C) = LS/(p{-}1)$, where $S = \sum_d n_d^2$. This is
the condition under which the alignment limit takes the
same value across all cosets. A strictly stronger
condition is the per-bin proportionality
$|C \cap B_d| = L \cdot n_d / (p{-}1)$ for every $d$,
which forces $T(C) = LS/(p{-}1)$ identically. In the
primitive root case ($L = p{-}1$), there is only one
coset and equidistributed bin sum is automatic.
\end{remark}

% ============================================================
\section{The Alignment Limit}

\begin{theorem}\label{thm:exact}
For any prime $p$ not dividing $b$ and any $b$-smooth
$m \ge 1$, the repetend alignment of $n = pm$ is
\[
\alpha(pm) = \frac{m - 1 + \dfrac{m \cdot T(C_m)}{L}}{pm - 1},
\]
where $L = \operatorname{ord}_p(b)$ and $C_m = u\langle b \rangle$
is the coset determined by $u = b^t/m$ with $m \mid b^t$.
\end{theorem}

\begin{proof}
The $m - 1$ multiples of $p$ in $\{1, \ldots, pm{-}1\}$ terminate,
each contributing alignment~$1$. The remaining $(p{-}1)m$ values
of $k$ are non-terminating. After the prefix of length $t$, the
repetend of $k/(pm)$ is the repetend of $(uk \bmod p)/p$, where
$u = b^t/m$.

The repetend of $1/(pm)$ has length~$L$. At the $i$-th
position, the digit is $\delta(r_i)$ where $r_i$ is the
$i$-th element of the coset $C_m$. A non-terminating
fraction $k/(pm)$ matches $1/(pm)$ at position~$i$ if and
only if $uk \bmod p$ lies in the same bin $B_{\delta(r_i)}$,
which contains $n_{\delta(r_i)}$ residues. Each residue
class mod~$p$ accounts for $m$ values of~$k$. Averaging the
match count over the $L$ positions:
\[
\text{aligned non-terminating} = m \cdot \frac{1}{L}
\sum_{r \in C_m} n_{\delta(r)} = \frac{m \cdot T(C_m)}{L}.
\]
Adding the $m - 1$ terminating fractions:
\[
\alpha(pm) \cdot (pm - 1)
= (m - 1) + \frac{m \cdot T(C_m)}{L}.
\qedhere
\]
\end{proof}

\begin{remark}
As $m \to \infty$, $\alpha(pm) \to (L + T(C_m))/(pL)$.
The limit depends on the coset $C_m$. When $b$ is a
primitive root modulo $p$, there is only one coset and
the limit is $(L + T)/(pL)$. When $L < p{-}1$, different
$b$-smooth values of $m$ may yield different cosets, and
the limit over all smooth $m$ does not exist in general.
\end{remark}

\begin{remark}
A concrete example: $p = 53$, $b = 10$, $L = 13$. There
are four cosets, paired by bin sum. The subgroup
$\langle 10 \rangle$ and one other coset have $T = 67$;
the cosets $2\langle 10 \rangle$ and $5\langle 10 \rangle$
have $T = 69$. The two values give two subsequential
limits, $80/689$ and $82/689$. For $m = 2^a$, $u = 5^a$,
and the coset alternates with $a \bmod 4$.
\end{remark}

\begin{corollary}\label{cor:closed-form}
When $b$ is a primitive root modulo $p$
($L = p{-}1$, one coset), the alignment limit exists
and equals
\[
\mathcal{L}(p,b) = \frac{p - 1 + S(p,b)}{p(p - 1)},
\]
where $S(p,b) = bq^2 + r(2q{+}1)$ with
$q = \lfloor(p{-}1)/b\rfloor$ and $r = (p{-}1) \bmod b$.
\end{corollary}

\begin{proof}
When $L = p{-}1$, there is only one coset and
$T = S$ (every residue is visited). Substituting into
Theorem~\ref{thm:exact} and taking $m \to \infty$:
$\mathcal{L} = (L + S)/(pL)
= (p{-}1+S)/(p(p{-}1))$.
\end{proof}

\begin{corollary}\label{cor:dp-limit}
For digit-partitioning primes ($p \le b{+}1$),
$S = p{-}1$ and $\mathcal{L}(p,b) = 2/p$.
\end{corollary}

\begin{proof}
By Corollary~\ref{cor:dp-recovery}, $S = p{-}1$.
Substituting into Corollary~\ref{cor:closed-form}:
$(p{-}1 + p{-}1)/(p(p{-}1)) = 2/p$.
\end{proof}

% ============================================================
\section{Enumeration in Base Ten}

\begin{table}[h]
\centering
\begin{tabular}{rrrcrl}
\toprule
$p$ & $L$ & $S$ & Equidist. & $\mathcal{L}(p)$ & Decimal \\
\midrule
 3 &  1 &   2 & \checkmark & $2/3$      & $0.6667$ \\
 7 &  6 &   6 & \checkmark & $2/7$      & $0.2857$ \\
11 &  2 &  10 & \checkmark & $2/11$     & $0.1818$ \\
13 &  6 &  16 & \checkmark & $7/39$     & $0.1795$ \\
17 & 16 &  28 & \checkmark & $11/68$    & $0.1618$ \\
19 & 18 &  34 & \checkmark & $26/171$   & $0.1520$ \\
23 & 22 &  50 & \checkmark & $36/253$   & $0.1423$ \\
29 & 28 &  80 & \checkmark & $27/203$   & $0.1330$ \\
31 & 15 &  90 & \checkmark & $4/31$     & $0.1290$ \\
37 &  3 & 132 & \checkmark & $14/111$   & $0.1261$ \\
41 &  5 & 160 & \checkmark & $5/41$     & $0.1220$ \\
43 & 21 & 186 & \checkmark & $110/903$  & $0.1218$ \\
47 & 46 & 216 & \checkmark & $130/1081$ & $0.1203$ \\
53 & 13 & 272 &            & $80/689$   & $0.1161$ \\
59 & 58 & 346 & \checkmark & $198/1711$ & $0.1157$ \\
61 & 60 & 366 & \checkmark & $7/61$     & $0.1148$ \\
\bottomrule
\end{tabular}
\caption{Alignment values for primes up to $61$ in
base~$10$. For primitive-root primes (checkmarked),
the limit is unique. For $p = 53$ ($L = 13$, four
cosets), the value shown is for the subgroup coset;
other cosets give different limits.}
\label{tab:limits}
\end{table}

% ============================================================
\section{The Golden Ratio Bound}

\begin{theorem}\label{thm:golden-bound}
For every prime $p \ge 5$ not dividing $b$, every base
$b \ge 2$, and every $b$-smooth $m \ge 1$:
\[
\alpha(pm) < \frac{1}{\varphi}.
\]
In particular, every subsequential alignment limit (over
$b$-smooth $m \to \infty$) is below $1/\varphi$.
\end{theorem}

\begin{proof}
For any coset $C$, the bin sum satisfies $T(C) \le L \cdot
\max_d n_d$ (each of the $L$ coset elements contributes
at most $\max_d n_d$). Since
$\max_d n_d \le \lceil(p{-}1)/b\rceil$, the alignment
at any $b$-smooth $m$ satisfies
\[
\alpha(pm)
\le \frac{m{-}1 + m c}{pm - 1}
\le \frac{1 + c}{p}
\le \frac{p+1}{2p},
\]
where $c = \lceil(p{-}1)/b\rceil$. The second inequality
holds because $(m{-}1+mc)(p) \le (1+c)(pm{-}1)$
simplifies to $p - 1 - c \ge 0$, which is true since
$c \le (p{-}1)/b \le p{-}1$.
For $p = 5$: $(p{+}1)/(2p) = 0.6 < 1/\varphi$.
For $p \ge 7$: $(p{+}1)/(2p) \le 4/7 < 1/\varphi$.
\end{proof}

Among primes $p \ge 3$, only $p = 3$ produces alignment
above $1/\varphi$. The bound holds for every coset and
every $b$-smooth $m$, without any equidistribution
hypothesis.

% ============================================================
\section{Remarks}

\subsection*{The role of equidistribution}

The equidistributed bin sum condition $T(C) = LS/(p{-}1)$
holds for every coset of $\langle b \rangle$ in
$(\mathbb{Z}/p\mathbb{Z})^{*}$ for all primes below $100$ in
base $10$ except $p = 53, 73, 79, 89$. These exceptions
arise when the orbit $\langle b \rangle$ is concentrated in
a subset of digit bins whose total size differs from the
average $LS/(p{-}1)$. Whether equidistribution holds for a
given $(p,b)$ is determined by the interaction between the
multiplicative structure of $(\mathbb{Z}/p\mathbb{Z})^{*}$
and the Euclidean partition induced by $\delta$;
see~\cite{iwaniec} for the general theory of character sums
and equidistribution in multiplicative groups. A
characterization of equidistributed pairs $(p,b)$ would
strengthen Corollary~\ref{cor:closed-form} into a theorem
with explicit domain.

\subsection*{Connection to the program}

The three papers together give:

\begin{itemize}
\item The golden ratio paper~\cite{paper1}: the golden ratio's
      minimal polynomial divides the self-referential cubic only at
      $p = 3$, producing the identity $\tau^4 = 2 - 3\tau$.
\item The digit-partitioning paper~\cite{paper2}: the alignment
      formula $(2m{-}1)/(pm{-}1)$ holds for all primes
      $p \le b{+}1$, characterized by injectivity of $\delta$.
\item The present paper: the alignment limit extends to all primes,
      with corrections from the squared bin sizes of $\delta$.
\end{itemize}

The next layer is the alignment of composite denominators $n$ with
multiple prime factors, where the Chinese Remainder Theorem governs
the interaction between prime contributions.

\subsection*{The bin-sum formula}

The quantity $S(p,b) = \sum n_d^2$ has a natural interpretation: it
is the expected number of remainders sharing a digit bin with a
uniformly random remainder. Equivalently, it measures the
\emph{collision probability} of the digit function. In Fourier-analytic terms, this sum
is the zeroth mode of a Dirichlet kernel
expansion~\cite{zygmund}. For
digit-partitioning primes, the collision probability is minimal
($S = p{-}1$, each remainder alone in its bin). As $p$ grows past
$b{+}1$, collisions increase quadratically with the bin sizes,
enriching the alignment limit.

% ============================================================
\begin{thebibliography}{9}

\bibitem{paper1}
A.~S.~Petty,
Why the golden ratio selects the prime three,
research note, 2020.

\bibitem{paper2}
A.~S.~Petty,
Digit-partitioning primes and the alignment formula,
research note, 2020.

\bibitem{davenport}
H.~Davenport,
\emph{Multiplicative Number Theory},
3rd ed., Springer, 2000.

\bibitem{zygmund}
A.~Zygmund,
\emph{Trigonometric Series}, vols.~I--II,
3rd ed., Cambridge University Press, 2002.

\bibitem{iwaniec}
H.~Iwaniec and E.~Kowalski,
\emph{Analytic Number Theory},
Amer.\ Math.\ Soc., 2004.

\end{thebibliography}

\end{document}
